\documentclass[11pt]{article}
\usepackage[T1]{fontenc}
\usepackage{lmodern}
\usepackage[margin=1in]{geometry}
\usepackage{amsmath,amssymb,amsthm}
\usepackage[hidelinks]{hyperref}
\hypersetup{pdftitle={Smooth nonquadratic objectives with periodically repeating Ritz steps},pdfauthor={},pdfsubject={AI-assisted research working paper, pending author review}}
\newtheorem{theorem}{Theorem}
\newtheorem{corollary}[theorem]{Corollary}
\newtheorem{lemma}[theorem]{Lemma}
\newcommand{\R}{\mathbb R}
\newcommand{\N}{\mathbb N}
\newcommand{\norm}[1]{\left\lVert#1\right\rVert}
\title{Smooth nonquadratic objectives with periodically repeating Ritz steps}
\author{}
\date{November 2026}
\begin{document}
\maketitle
\begin{center}\small\emph{AI-assisted working paper. Author verification pending.}\end{center}
\begin{abstract}
A delayed Ritz gradient method on a positive definite quadratic can
repeat its step sizes even when its gradient direction does not return
after one period. Suppose the active spectral components have the same absolute
multiplier over one period. The iterates then lie on finitely many
lines. We construct a smooth perturbation that vanishes along those
lines with all its derivatives, while remaining nonzero arbitrarily
close to the minimizer. Adding it produces a globally strongly convex
nonquadratic objective. For a specified compatible history, a method
that uses only data at visited points then generates the same stored
gradients, step sizes, and iterates. We apply the construction to a six-sweep orbit of
limited-memory steepest descent with memory two and preserve its local
lower bound on the gradient decay factor per scalar step.
\end{abstract}

\section{A periodic orbit supplies a finite set of lines}

Consider \(f_A(x)=\tfrac12 x^\top A x\) with \(A\) real symmetric and
positive definite on \(\R^d\), \(d\ge 2\). We number scalar gradient steps
starting at the beginning of a specified compatible prefix. Let
\(g_k=Ax_k\) and \(x_{k+1}=x_k-\alpha_k g_k\), where \(\alpha_k>0\).
In memory-two limited-memory steepest descent (LMSD), two step sizes
are computed from the preceding block of two scalar updates. Each
block of two updates is called a sweep. The steps are the reciprocals
of the two Ritz values obtained by projecting the Hessian onto the
span of the two gradients used in that preceding block. In the six-sweep example, the prefix consists of the first two
steps. At its end, the method has three iterates, their gradients,
and the two reciprocal Ritz values for the next sweep. By a compatible prefix
we mean that the stored samples were produced by the displayed scalar
steps and that the next Ritz values were computed by the prescribed
history rule. Any data from before the displayed prefix that are needed to
start the method are part of the specified initialization.

Assume \(g_0\ne0\) and the step sizes repeat with integer period \(T\ge1\)
along the quadratic
trajectory. Define the period polynomial by
\[
 q(t)=\prod_{k=0}^{T-1}(1-\alpha_k t).
\]
Suppose each eigenvalue \(\lambda\) with a nonzero component in \(g_0\)
satisfies \(|q(\lambda)|=G\), for a common number \(0<G<1\). Denote by
\(E\) the sum of the corresponding eigenspaces. Each \(g_k\) and \(x_k\)
belongs to \(E\). The restriction \(S=q(A)|_E/G\) is symmetric,
commutes with \(A|_E\), and satisfies \(S^2=I_E\). In particular
it is an isometry on \(E\).

For \(0\le r<T\), take the first \(r\) scalar factors in the period.
They commute with \(q(A)\), so for \(n\in\N=\{0,1,\ldots\}\)
\begin{equation}
 g_{nT+r}=G^n S^n g_r,\qquad
 x_{nT+r}=G^n S^n x_r .
 \label{eq:period}
\end{equation}
An active component can never be killed within the period. Killing it
would give \(q(\lambda)=0\), contrary to \(G>0\). Thus \(x_r\ne0\), and
every \(x_k\) belongs to one of the at most \(2T\) lines
\begin{equation}
 \mathcal L=\{\operatorname{span}(x_r),\operatorname{span}(Sx_r)
                 :0\le r<T\}.
 \label{eq:lines}
\end{equation}
The signs in \(S\) explain why both lines per phase may be needed.
The formula also gives a uniform bound on the entire set of iterates
and on its positive rescalings with scale at most one.

\section{A smooth perturbation that preserves the visited data}

\begin{lemma}\label{lem:flat}
Given a finite nonempty collection of lines through the origin in
\(\R^d\), \(d\ge2\), there is a nonnegative smooth function \(\psi\)
of compact support whose derivatives of every order vanish on the
union of the lines. The function is positive at some point in every
neighborhood of the origin.
\end{lemma}

\begin{proof}
Let \(R_L\) be the orthogonal projection onto the orthogonal complement
of \(L\). Define \(b(0)=0\) and \(b(t)=\exp(-1/t)\) for \(t>0\).
Extend \(b\) by zero to \(t<0\). For \(t>0\), every derivative of \(b\)
is a finite sum of multiples of \(t^{-j}\exp(-1/t)\), as follows by
differentiating and increasing \(j\). Each such term tends to zero
as \(t\downarrow0\), even after division by any positive power of \(t\).
Induction using the difference quotient at zero proves that the
extension is smooth with \(b^{(k)}(0)=0\) for every \(k\ge0\).

The function \(x\mapsto b(\norm{R_Lx}^2)\) is smooth. Its value is
zero on \(L\), and every term in its differentiated chain rule has a
factor \(b^{(j)}(0)\) there. Let \(\chi\) be a smooth nonnegative
compactly supported function equal to one near zero and put
\[
 \psi(x)=\chi(x)\prod_{L\in\mathcal L} b(\norm{R_Lx}^2).
\]
The product rule retains the vanishing derivatives on each \(L\).
A finite union of proper linear subspaces has empty interior. For
example, each subspace has zero \(d\)-dimensional Lebesgue measure,
so the union cannot contain a ball. In every sufficiently small
ball there is a point off all lines where all factors are positive.
\end{proof}

Only agreement of the function, gradient, and Hessian is needed for
the basic Ritz application. For agreement through order two a
polynomial factor \(\norm{R_Lx}^4\) suffices. Writing \(v=R_Lx\), its
gradient is \(4\norm{v}^2v\) and its Hessian is
\(8vv^\top+4\norm{v}^2R_L\). The smooth factor used above also covers a
rule that reads a finite number of higher derivatives at visited
points.

\begin{theorem}\label{thm:lift}
Assume the quadratic orbit and its period satisfy the conditions
of Section~1. Fix a deterministic exact-arithmetic implementation
that uses only derivatives and values at generated iterates. Suppose a
fixed finite initialization specifies its starting phase, and
rescaling that initialization on the quadratic by \(s>0\) gives
the trajectory \(s x_k\) with the same sequence of step sizes.
Any further observations used to produce the starting phase must
come from finitely many points whose locations scale by \(s\). There is a globally strongly convex
\(F\in C^\infty(\R^d)\) with unique minimizer \(0\) and
\(\nabla^2 F(0)=A\) that agrees with \(f_A\) to every derivative order
at \(0\) and is nonquadratic on every neighborhood of \(0\). For
every \(s>0\), this implementation on \(F\), with the corresponding
scaled matching prefix and the same deterministic conventions,
reproduces \(s x_k\) at every subsequent step.
\end{theorem}

\begin{proof}
Add to \(\mathcal L\) in \eqref{eq:lines} the lines through any
nonzero points used in the fixed finite initialization. This is
still a finite set, and it contains all their positive rescalings.
Use Lemma~\ref{lem:flat} with that set, and let
\(M=\sup_x\norm{\nabla^2\psi(x)}_2\). Compact support and smoothness
give \(M<\infty\). If \(M=0\), then \(\psi\) is affine, which together
with compact support makes it zero. Choose \(0<\eta<\lambda_{\min}(A)/M\)
and put \(F=f_A+\eta\psi\). For each \(z\in\R^d\),
\[
 z^\top\nabla^2F(x)z
 \ge (\lambda_{\min}(A)-\eta M)\norm{z}^2 .
\]
This gives global strong convexity. The gradient at zero vanishes
and every derivative at zero matches the quadratic. If \(F\) were
quadratic on a neighborhood of zero, its value, gradient and Hessian
there would determine that polynomial as \(f_A\), contradicting
Lemma~\ref{lem:flat}.

Fix \(s>0\). Every point \(sx_k\) belongs to the same set of lines.
By scale equivariance, running on \(f_A\) from the scaled quadratic
prefix produces \(sx_k\). Every observation on that prefix is equal
for \(F\) and \(f_A\), which initializes the same stored state,
including the steps already computed for the next sweep. If the
two states agree before a step, the same rule computes the same
step size from the same inputs. Both move to the same next point
using the same gradient. The new observation again agrees, so the
stored states after that step agree. This proves the claim by
induction.
\end{proof}

The data condition in the theorem includes pure, unweighted
memory-two LMSD with a deterministic order of Ritz values and an
exact rank rule. Equality of the stored iterates and gradients makes
the gradient differences, projected pencil, rank decision and root
order identical. Scale equivariance in the quadratic follows because
scaling all history vectors preserves their span. An absolute
stopping tolerance, rejected trial evaluations, or other off-orbit
queries can change the run and need separate analysis.

\section{The six-sweep example for memory two}

Take a sufficiently small \(\varepsilon>0\) and set
\[
 A_\varepsilon=\operatorname{diag}(1/2,\ 1/2+\varepsilon,\
                                  1-\varepsilon,\ 1).
\]
The six-sweep construction of \cite[Theorem 8.1 and Appendix R]{sharp} gives a full-support quadratic orbit and a
compatible entry \(u_1=P_5(A_\varepsilon)u_0\).
The six sweep polynomials are \(P_0,\ldots,P_5\), with each
\(P_i(t)=(1-t/\vartheta_{i,1})(1-t/\vartheta_{i,2})\).
The displayed prefix starts at \(g_0=u_0\). Apply the steps
\(1/\vartheta_{5,1}\) and \(1/\vartheta_{5,2}\) in that order to
obtain \(g_1\) and \(g_2=u_1\). To see what the pure
method observes, let \(B=[g_0,g_1]\), and let \(\alpha_0,\alpha_1\)
be the two prefix steps. The two columns of \(AB\) are exactly
\((g_0-g_1)/\alpha_0\) and \((g_1-g_2)/\alpha_1\). If \(B=VR\)
is a reduced QR factorization of this rank-two history, the symmetric
matrix whose eigenvalues yield the next steps is
\(V^\top(AB)R^{-1}=V^\top A V\). The three sampled gradients and the
two steps therefore determine this projected problem. It produces
\(P_0\) on the prescribed branch. The ordered sequence from the start
of this prefix uses \(P_5\)
followed by \(P_0,\ldots,P_4\), and repeats. The choice of order
within each sweep affects intermediate points but leaves the
polynomial product unchanged.

Write \(Q=\prod_{i=0}^{5}P_i\). If \(g\) is the gradient
at the start of a period, its spectral weights are
\(w_i=|\langle g,e_i\rangle|^2/\norm{g}^2>0\), where \(e_i\)
is the \(i\)th coordinate eigenvector. Applying \(Q(A_\varepsilon)\) changes
the weights to
\[
 \frac{w_i Q(\lambda_i)^2}{G^2},\qquad
 G^2=\sum_j w_j Q(\lambda_j)^2 .
\]
The weights return to their original values by the six-cycle
construction. Cancelling \(w_i\) therefore gives
\(Q(\lambda_i)^2=G^2\) at all four nodes. The six-sweep construction gives
\begin{equation}
 G^{1/12}
   = 2^{-1/2}\varepsilon^{3/8}
       \bigl(1+O(\sqrt\varepsilon)\bigr).
 \label{eq:factor}
\end{equation}
Here is the initialization and time count used below. Choose
\(u_0\) with the four normalized squared components prescribed by
the phase-zero solution in \cite[Appendix R]{sharp}. In the diagonal
model one can take the positive square roots as its coordinates.
Set \(x_0=A_\varepsilon^{-1}u_0\) and take the two ordered reciprocal
Ritz steps from \(P_5\). The pair \((u_0,u_1)\), together with these
two scalar steps, is the compatible entry. The three measured
gradients now determine the two Ritz values defining \(P_0\) as above.
Start the lower-bound clock at \(g_2=u_1\), so its time
\(\ell\) corresponds to \(k=\ell+2\) on the clock of
\eqref{eq:period}. The next twelve scalar updates consume
\(P_0,\ldots,P_5\) in this order. The six-cycle construction gives all six of these
polynomials by recomputing the history at each boundary, not by
preassigning the subsequent updates to the nonlinear run.
For each \(\varepsilon\), Theorem~\ref{thm:lift} keeps that rule and
its entire trajectory on one nonlinear objective at every positive
scale.

\begin{corollary}\label{cor:local}
Fix that nonlinear objective and any neighborhood \(U\) of zero.
Let \(H_-=(h_0,h_1,h_2)\) be the three stored gradients at the end
of a compatible prefix, and let \(\widehat x_\ell\) denote the
iterates counted from that phase. For \(0<\gamma<G^{1/12}\), no
finite \(C\), independent of the phase state and \(\ell\), makes
\[
 \norm{\nabla F(\widehat x_\ell)}
 \ \le\ C\gamma^\ell \max_{0\le j\le2}\norm{h_j}
 \qquad(\ell\ge0)
\]
hold for all such initialized trajectories whose stored points
and future iterates lie in \(U\).
\end{corollary}

\begin{proof}
The unscaled iterates are bounded by \eqref{eq:period}. Hence a
single small \(s>0\) places the entire run and the three stored
points in \(U\). Set \(H_-^{(s)}=(s g_0,s g_1,s g_2)\). At time
\(\ell=12n\), the isometry of \(S\) and \eqref{eq:period} give
\(\norm{\nabla F(\widehat x_{12n})}=s G^n \norm{g_2}\).
The proposed estimate would force
\[
 C\ \ge\
 \frac{\norm{g_2}}{\max_{0\le j\le2}\norm{g_j}}
 \left(\frac{G}{\gamma^{12}}\right)^n
\]
for every \(n\). The prefactor is positive and the base exceeds
one. This contradiction already holds for the single scaled phase
state selected above.
\end{proof}

\section{A real-analytic variant for finitely many derivatives}

The orbit can also be preserved by a real-analytic objective if
the implementation requires agreement only through a fixed derivative
order \(m\). Take an integer \(k\) with \(2k>\max(m,2)\).
In the proof of
Theorem~\ref{thm:lift}, one may use
\[
 \psi_{\mathrm{an}}(x)
 = e^{-\norm{x}^2}\prod_{L\in\mathcal L}\norm{R_Lx}^{2k}.
\]
To verify the required vanishing, expand the \(L\)-factor at a point
\(x_0\in L\). It is the homogeneous polynomial
\(\norm{R_L(x-x_0)}^{2k}\) in the increment, so every derivative of
order below \(2k\) is zero there. The product rule proves the same
vanishing for \(\psi_{\mathrm{an}}\). Every second derivative is a
polynomial times \(e^{-\norm{x}^2}\), hence bounded on \(\R^d\).
The Hessian estimate in Theorem~\ref{thm:lift} now gives a globally
strongly convex, real-analytic \(F\) that is nonquadratic in every
neighborhood of zero and preserves the same trajectory. For the
plain first- and second-derivative version of the method, \(k=2\)
suffices. The construction from Lemma~\ref{lem:flat} is needed for
agreement through every derivative order.

\section{Scope and relation to earlier work}

The lower bound concerns a prescribed compatible history and a
pure visited-point implementation. A procedure that starts from a
single point with two ordinary steepest-descent line searches may
generate a different history. The argument establishes neither
a basin of initial histories nor persistence under perturbation of
step sizes, function evaluations, or roundoff. Equation
\eqref{eq:factor} concerns the specified six-cycle. It gives a
lower witness and leaves the nonlinear worst-case upper rate open.

Ritz sweeps for nonquadratic problems are studied in
\cite{fletcher,ferrandi}. The pure realization used here omits the
line-search and safeguard mechanisms used in implementations of such
methods. Smooth strongly convex interpolation \cite{taylor} supplies
finite-data constructions for performance estimation. Our
perturbation prescribes data on the finitely many lines that contain
an infinite trajectory and all its positive rescalings. The result
depends on the complete delayed-state return proved for the chosen
quadratic branch. Gu \cite{gu} studies convergence under cyclic
reuse of historical Ritz values on quadratics. The orbit used here
repeats the sweep \emph{pattern} after six sweeps of the usual
memory-two rule, which consumes each computed Ritz pair once.

The general lifting criterion can be applied to another periodic
quadratic history once its full return, finite matching initialization,
and scalar step count have been proved for that history. The quadratic six-cycle is the one proved in \cite{sharp}.
The associated nonlinear six-cycle with matching derivatives
through order two was derived in an earlier unpublished version
of this work. This paper presents that derivation in a
cycle-independent form and adds perturbations that either match
derivatives of all orders or are real-analytic and match any
prescribed finite order.

\paragraph{Status and AI disclosure.}
This is an AI-assisted research working paper based on an existing
research derivation. It is pending human author review. The companion
formalization checks state replay and the discrete return recursion
under explicit hypotheses. The analytic construction, the quadratic
orbit, and its asymptotic rate are proved or cited in the manuscript
and are outside that machine-checked scope.

\begin{thebibliography}{9}
\bibitem{sharp}
\emph{Sharp spectral bounds and regeneration in limited-memory steepest
descent}, anonymized research working paper, November 2026.
\href{https://taitai2025.github.io/assets/documents/research/lmsd-spectral.pdf}
{Archived version with proof appendices}.
\bibitem{fletcher}
R. Fletcher,
A limited memory steepest descent method,
\emph{Mathematical Programming}, 135 (2012), pp.\ 413--436,
\href{https://doi.org/10.1007/s10107-011-0479-6}{doi:10.1007/s10107-011-0479-6}.
\bibitem{ferrandi}
G. Ferrandi and M.~E. Hochstenbach,
Limited memory gradient methods for unconstrained optimization,
\href{https://arxiv.org/abs/2308.15145}{arXiv:2308.15145}, 2023.
\bibitem{taylor}
A.~B. Taylor, J.~M. Hendrickx, and F. Glineur,
Smooth strongly convex interpolation and exact worst-case performance
of first-order methods,
\href{https://arxiv.org/abs/1502.05666}{arXiv:1502.05666}.
\bibitem{gu}
R. Gu, Superlinear convergence of cyclic gradient methods with
historical Ritz step sizes,
\href{https://arxiv.org/abs/2610.32705v2}{arXiv:2610.32705v2}, 2026.
\end{thebibliography}
\end{document}
